% Generated draft: requires mathematical and source review.
\documentclass[11pt]{article}
\usepackage[margin=25mm]{geometry}
\usepackage{fontspec}
\setmainfont{Latin Modern Roman}
\usepackage{amsmath,amssymb,mathtools,unicode-math}
\setmathfont{Latin Modern Math}
\usepackage{enumitem,needspace,xurl,hyperref}
\setlength{\parindent}{0pt}
\setlength{\parskip}{5pt}
\setlength{\emergencystretch}{3em}
\newcommand{\N}{\mathbb{N}}
\newcommand{\Z}{\mathbb{Z}}
\newcommand{\Q}{\mathbb{Q}}
\newcommand{\R}{\mathbb{R}}
\newcommand{\C}{\mathbb{C}}
\newcommand{\F}{\mathbb{F}}
\newcommand{\A}{\mathbb{A}}
\newcommand{\PP}{\mathbb{P}}
\newcommand{\E}{\mathbb{E}}
\newcommand{\eps}{\varepsilon}
\newcommand{\dd}{\,\mathrm{d}}
\begin{document}
\section*{UnsolvedMath ID: 10900013}
\textbf{Source problem code:} catalogue entry
\subsection{Definitions and mathematical statement}
Let \(M=\mathbb H^3/\Gamma\) be a cusped finite-volume hyperbolic \(3\)-manifold. A closed quasi-Fuchsian surface is understood as a closed surface immersed \(\pi_1\)-injectively in \(M\), whose surface subgroup \(H\leq\Gamma\) is quasi-Fuchsian; its limit set \(\Lambda H\subset\partial_\infty\mathbb H^3\cong S^2\) is a Jordan curve. Lifts correspond to conjugates, with limit sets \(g\Lambda H\) for \(g\in\Gamma\).

Agol asks whether such surfaces cubulate \(M\) except for its cusps. The source's stated equivalent geometric formulation is
\[
 \forall p\ne q\in S^2\quad\exists H\text{ from a closed quasi-Fuchsian surface}
 \quad\exists g\in\Gamma:
 \ p,q\text{ lie in different components of }S^2\setminus g\Lambda H.
\]
In particular both points must avoid that separating limit set. The points must be distinct for separation to be possible; this implicit qualification is made explicit.

The source says a positive answer would give a cocompact action of \(\Gamma\) on a CAT\((0)\) cube complex whose only point stabilizers are parabolic subgroups. Here a cusp subgroup fixes a parabolic point at infinity; the precise meaning of permitted stabilizers, including trivial ones and subgroups versus full cusp groups, is not spelled out. The statement also does not specify whether finitely many immersed surfaces are required to furnish all separating lifts, whether an orientable/torsion-free convention is assumed throughout, or the exact meaning of cubulating except for cusps. The claimed equivalence and consequence are retained as source assertions, without replacing this property by generic virtual specialness or by existence of a single surface.
\subsection{Short English statement}
Do closed quasi-Fuchsian surfaces in every cusped finite-volume hyperbolic \(3\)-manifold separate every distinct pair at infinity through suitable lifts, yielding the source's cubulation with parabolic stabilizers?
\subsection{Sources}
\begin{itemize}
\item Agol, attributed item 3.5 in Delp, Hoffoss and Manning, \emph{Problems In Groups, Geometry, and Three-Manifolds} (2015), supplied PDF page 2. \url{https://arxiv.org/pdf/1512.04620#page=2}; status link \url{https://arxiv.org/abs/1512.04620}. Recorded access 2026-07-29; extraction source-tex checked against rendered PDF text, as reported in the record.
\item The source points to Masters--Zhang and Baker--Cooper surface constructions. The supplied summary additionally mentions Agol--Wise virtual specialness and Kahn--Markovic-type work, without complete bibliography or an exact theorem resolving the stated separation property.
\item Local UnsolvedMath registry, \texttt{data/unsolvedmath/problems.json}, permanent ID \(10900013\), exact source code \texttt{catalogue entry}; supplied partially\_solved/PARTIAL-PROGRESS status.
\end{itemize}
\end{document}
